% ===================================================================
%  TABELUL Nr. 4 — Maturitatea și Bătrânețea.
%  Traduction 2026 d'apres texto/io, controlee sur texto/fr, texto/en
%  et texto/es. Consigne : voir texto/ro/10-tabelul-01.tex.
%
%  LE TABLEAU DE LA PARENTE, ET LE ROUMAIN Y EST PLUS PRECIS QUE LE
%  FRANCAIS : « cumnat » le beau-frere, « cuscru » le pere du gendre,
%  « socru » et « soacră » les beaux-parents, « nora » la belle-fille,
%  « ginere » le gendre. Six degres, six mots simples, la ou le
%  francais compose tout avec « beau- ».
% ===================================================================

%%K t04-apar-1 apar
\VUsaut{4.0mm}
\VUtitre{15.0pt}{60}{TABELUL Nr. 4}
\VUsaut{1.6mm}
\VUtitre{11.4pt}{0}{Maturitatea și Bătrânețea.}
\VUsaut{3.2mm}

%%K t04-tit-1 sub
\VUcentre{11.4pt}{0}{\textit{Prima scenă.}}
\VUsaut{1.6mm}
\VUtitre{11.4pt}{0}{Ospățul de nuntă.}
\VUsaut{2.6mm}

%%K t04-tit-2 sub
\VUpk{10.2pt}{40}{Toamna.}
\VUsaut{3.0mm}

%%K t04-c1-01-1 p
1. --- Tinerii au crescut. Unul dintre ei, Ioan, se însoară. Suntem la
\VUgras{ospățul de nuntă}, care a strâns la aceeași masă familiile mirelui și ale
miresei.

%%K t04-c1-02-1 p
2. --- Este \VUgras{toamnă}, luna \VUgras{septembrie}. Vântul rece al lui
\VUgras{octombrie} și \VUgras{noiembrie} nu a smuls încă de pe câmp frunzișul
verde. Aerul serii este blând și înmiresmat; de aceea ospățul se dă sub
\VUgras{verandă} \textsuperscript{(8)}, care dă spre grădină.

%%K t04-c1-03-1 p
3. --- În depărtare, pe \VUgras{deal} \textsuperscript{(3)}, se află casa
părinților lui Ioan, un \VUgras{conac} elegant \textsuperscript{(1)}, ascuns în
verdeață; vii frumoase, care dau un vin ales, acoperă coasta dealului. Viticultorul
le îngrijește și le răsfață.

%%K t04-c1-04-1 p
4. --- Deasupra vițelor trece în zbor un cârd de \VUgras{potârnichi}
\textsuperscript{(2)}, speriate de apropierea \VUgras{pădurarului}
\textsuperscript{(5)}. Iar acesta, cu \VUgras{pușca} la subsuoară și cu
\VUgras{tolba} pe umăr, scotocește \VUgras{tufișurile} \textsuperscript{(4)} cu
\VUgras{câinele} său \textsuperscript{(6)}. Păzește moșia și nu îi lasă pe
braconieri să nimicească \VUgras{vânatul}.

%%K t04-c1-05-1 p
5. --- Pe dinafară se cațără pe \VUgras{verandă} o \VUgras{viță de boltă}, din
care atârnă \VUgras{ciorchini} grei \textsuperscript{(7)}.

%%K t04-c1-06-1 p
6. --- Florile frumoase de toamnă care împodobesc \VUgras{masa}
\textsuperscript{(16)} sunt \VUgras{crizanteme} \textsuperscript{(9)};
\VUgras{tatăl} miresei \textsuperscript{(10)} și-a pus una la butoniera fracului.

%%K t04-c1-07-1 p
7. --- Ospățul se apropie de sfârșit. \VUgras{Servitorul}
\textsuperscript{(11)} începe să aducă farfuriile pentru desert și va strânge
îndată părțile de tacâm care nu mai sunt de trebuință --- \VUgras{lingurile}
\textsuperscript{(14)}, \VUgras{furculițele} \textsuperscript{(12)},
\VUgras{cuțitele} \textsuperscript{(13)}, \VUgras{platourile}
\textsuperscript{(15)}.

%%K t04-tit-3 sub
\VUpk{10.2pt}{40}{Familia.}
\VUsaut{3.0mm}

%%K t04-c1-08-1 p
8. --- Toți au înfățișarea veselă, iar zâmbetul cu care \VUgras{mama}
\textsuperscript{(17)} mirelui își privește copiii arată cât este de fericită.

%%K t04-c1-09-1 p
9. --- Această căsătorie unește două familii bogate din ținut. Ioan,
\VUgras{soțul} \textsuperscript{(18)}, este \VUgras{fiul} unui mare proprietar de
pământ și ajunge \VUgras{ginerele} unui fabricant înstărit.

%%K t04-c1-10-1 p
10. --- \VUgras{Soții} sunt tineri, și totuși au fost multă vreme
\VUgras{logodiți}. \VUgras{Tânăra soție} \textsuperscript{(19)}, Marta, este
drăgălașă în \VUgras{rochia} ei \VUgras{albă}, împodobită cu flori de portocal.

%%K t04-c1-11-1 p
11. --- \VUgras{Socrul} ei \textsuperscript{(20)} este bucuros de o \VUgras{noră}
atât de bună, iar \VUgras{soacra} lui Ioan \textsuperscript{(21)} zâmbește
văzându-și \VUgras{fiica} atât de frumoasă.

%%K t04-c1-12-1 p
12. --- La capătul mesei stau ceilalți \VUgras{oaspeți} \textsuperscript{(22)}:
\VUgras{rude, prieteni, veri}. \VUgras{Șeful de sală} \textsuperscript{(23)}, cu
șervetul pe braț, îi servește cu sprinteneală.

%%K t04-tit-4 sub
\VUpk{10.2pt}{40}{Prietenii.}
\VUsaut{3.0mm}

%%K t04-c1-13-1 p
13. --- În dreapta mesei se află o \VUgras{domnișoară} \textsuperscript{(24)},
\VUgras{prietena} miresei, apoi fratele miresei, \VUgras{cavalerul} ei de onoare
\textsuperscript{(25)}. Stă de vorbă cu \VUgras{domnișoara} sa de onoare
\textsuperscript{(26)}, sora mirelui, care prin această căsătorie ajunge
\VUgras{cumnata} Martei. Este o tânără fermecătoare. Are \VUgras{bijuterii}
frumoase, un \VUgras{colier de perle} și o \VUgras{brățară} de aur.

%%K t04-c1-14-1 p
14. --- Lângă ei domnul care s-a ridicat și a înălțat paharul este
\VUgras{prietenul} \textsuperscript{(27)} familiei. Este unul dintre
\VUgras{martorii mirelui}. Rostește un \VUgras{toast}, bea pentru
\VUgras{sănătatea} tinerilor și le urează fericire îndelungată. Este în ținută de
gală: frac și \VUgras{pantofi de lac} \textsuperscript{(28)}. O duce de braț pe
\VUgras{bunica} \textsuperscript{(29)}, care este \VUgras{văduvă}.

%%K t04-c1-15-1 p
15. --- Tânărul în \VUgras{frac} \textsuperscript{(31)}, cu mustăți negre subțiri,
este \VUgras{cumnatul} \textsuperscript{(30)} lui Ioan. Este un inginer de seamă.
Stă lângă o \VUgras{doamnă} foarte elegantă \textsuperscript{(33)},
\VUgras{sora mai mare} a Martei și mama fetiței drăgălașe care vine de braț cu
vărul ei să întindă o floare și să \VUgras{ureze} \VUgras{bună seara}. Această
doamnă are \VUgras{cercei} frumoși și o \VUgras{coafură} elegantă.

%%K t04-c1-16-1 p
16. --- Micul văr, atât de caraghios în \VUgras{bluza} lui \textsuperscript{(36)}
și în \VUgras{pantalonașii} largi \textsuperscript{(37)}, este vrednic de milă.
Este \VUgras{orfan} \textsuperscript{(35)}. Și-a pierdut tatăl și mama cu totul
mic. Unchiul lui îi este \VUgras{tutore} \textsuperscript{(38)} și a rămas
neînsurat ca să crească bietul copil. Este un om bun. Este destul de pleșuv și
foarte miop; poartă \VUgras{pince-nez}.

%%K t04-tit-5 sub
\VUsaut{2.4mm}
\VUpk{10.2pt}{40}{Desertul.}
\VUsaut{3.0mm}

%%K t04-c1-17-1 p
17. --- În spatele oaspeților, pe o masă acoperită cu o \VUgras{față de masă},
este așezat \VUgras{desertul} \textsuperscript{(39)}. Într-o vază mare fructe:
\VUgras{piersici} \textsuperscript{(40)}, \VUgras{pere} \textsuperscript{(41)},
\VUgras{mere} \textsuperscript{(42)}, \VUgras{caise} \textsuperscript{(43)} și
\VUgras{struguri} \textsuperscript{(44)}. Alături --- \VUgras{ananași}
\textsuperscript{(45)}, un \VUgras{tort} minunat \textsuperscript{(46)},
\VUgras{sticle de lichior} \textsuperscript{(48)} și de \VUgras{șampanie}
\textsuperscript{(49)}, precum și teancuri de \VUgras{farfurii} curate
\textsuperscript{(47)}.

%%K t04-c1-18-1 p
18. --- \VUgras{Paharnicul} \textsuperscript{(50)} destupă o \VUgras{sticlă}
\textsuperscript{(52)} de vin vechi cu \VUgras{tirbușonul}
\textsuperscript{(51)}. \VUgras{Dopul} \textsuperscript{(53)} iese, pare-se, foarte
greu. Acest paharnic are pe braț un \VUgras{șervet} \textsuperscript{(54)}.

%%K t04-c1-19-1 p
19. --- După ospăț bărbații vor trece în salonul de fumat, iar doamnele se vor duce
să se îmbrace pentru bal.

%%K t04-tit-6 sub
\VUsaut{5.0mm}
\VUcentre{11.4pt}{0}{\textit{A doua scenă.}}
\VUsaut{1.6mm}
\VUpk{11.4pt}{40}{Ziua onomastică a bunicului.}
\VUsaut{2.6mm}

%%K t04-tit-7 sub
\VUpk{10.2pt}{40}{Vizita.}
\VUsaut{3.0mm}

%%K t04-c2-01-1 p
1. --- Au trecut zece ani. Părinții lui Ioan au îmbătrânit și își iubesc cu
duioșie cei trei nepoți: doi \VUgras{nepoți} \textsuperscript{(2)} și o
\VUgras{nepoată} \textsuperscript{(3)}. Fiindcă astăzi este \VUgras{ziua
onomastică a bunicului}, copiii au venit să îl felicite.

%%K t04-c2-02-1 p
2. --- Micul Petru este încă foarte mic; are numai trei ani. Vede cum vine spre el
\VUgras{pisica} \textsuperscript{(1)}. Ar vrea să îi mângâie \VUgras{blana}
mătăsoasă și frumoasă și \VUgras{coada} lungă; nici nu îi trece prin cap că sub
aceste \VUgras{lăbuțe} elegante sunt gheare ascuțite și rele, care îl pot zgâria.

%%K t04-c2-03-1 p
3. --- Sora lui, Gabriela, care îl ține de mână, este mult mai serioasă, iar
Marcel, în \VUgras{paltonul} lui mare \textsuperscript{(4)} și în \VUgras{cureaua}
de piele \textsuperscript{(5)}, pare cu totul mare, cu toți pantalonii scurți din
care i se văd \VUgras{ciorapii} \textsuperscript{(6)}.

%%K t04-c2-04-1 p
4. --- Tatăl lor și-a pus \VUgras{paltonul} gros de iarnă \textsuperscript{(7)} cu
\VUgras{guler de blană} \textsuperscript{(8)}, iar în \VUgras{buzunar}
\textsuperscript{(9)} are un fular. Soția lui are o pălărie de pâslă cu
\VUgras{pene} \textsuperscript{(10)}, o \VUgras{jachetă} \textsuperscript{(11)}, un
\VUgras{boa} de blană \textsuperscript{(12)} și un \VUgras{manșon}
\textsuperscript{(13)}.

%%K t04-c2-05-1 p
5. --- Este foarte frig, căci a venit iarna. \VUgras{Zăpada}
\textsuperscript{(19)} a căzut toată ziua, iar acoperișurile caselor sunt cu totul
albe. Odată cu \VUgras{noaptea} gerul se întărește și mai tare. Pe cer
strălucesc \VUgras{luna} \textsuperscript{(17)} și \VUgras{stelele}
\textsuperscript{(18)}, străpungând \VUgras{întunericul}.

%%K t04-tit-8 sub
\VUsaut{4.2mm}
\VUpk{10.2pt}{40}{Boala.}
\VUsaut{3.0mm}

%%K t04-c2-06-1 p
6. --- Bietul \VUgras{orb} \textsuperscript{(20)}, care trece piața dinaintea
\VUgras{farmaciei} \textsuperscript{(16)}, călăuzit de \VUgras{pudelul} său
\textsuperscript{(21)}, este foarte nenorocit. Justina, \VUgras{slujnica}
\textsuperscript{(14)}, aduce din bucătărie o \VUgras{ceașcă}
\textsuperscript{(15)} de \VUgras{ceai de plante} foarte fierbinte, îndulcit din
belșug.

%%K t04-c2-07-1 p
7. --- \VUgras{Bunica} \textsuperscript{(23)}, care vrea să ia de pe masă
\VUgras{ochelarii} săi \textsuperscript{(26)} ca să citească \VUgras{rețeta}
\textsuperscript{(27)} scrisă adineauri de \VUgras{doctor}
\textsuperscript{(25)}, se ridică din fotoliu, sprijinindu-se în \VUgras{bastonul}
său \textsuperscript{(24)}.

%%K t04-c2-08-1 p
8. --- Sufere adesea de mici \VUgras{beteșuguri}, de \VUgras{amețeli}, de
\VUgras{răceli} și de \VUgras{dureri de dinți}; picioarele îi sunt slabe, iar
ochii nu mai sunt buni. Și totuși este bucuroasă să îl necăjească pe bătrânul ei
\VUgras{doctor}, care vrea mereu să îi prescrie o \VUgras{doctorie}
\textsuperscript{(28)}. Astăzi este foarte fericită că are în jur rudele tinere.

%%K t04-c2-09-1 p
9. --- În odaia bine închisă este plăcut. În \VUgras{căminul} de marmură
\textsuperscript{(29)} arde un \VUgras{foc bun} \textsuperscript{(30)}, și este
bine la suflet în \VUgras{lumina} blândă a \VUgras{lămpii} abia aprinse
\textsuperscript{(31)}.

%%K t04-tit-9 sub
\VUsaut{4.2mm}
\VUpk{10.2pt}{40}{Bunicul.}
\VUsaut{3.0mm}

%%K t04-c2-10-1 p
10. --- Până și \VUgras{bunicul} \textsuperscript{(33)} a uitat că a fost bolnav,
că \VUgras{reumatismul} îl ține țintuit în \VUgras{fotoliu}
\textsuperscript{(34)} și că chiar ieri a avut \VUgras{friguri și leșinuri}. Nu
mai ține minte că iarna trecută a avut \VUgras{aprindere de plămâni} și că o
\VUgras{tuse} crâncenă și chinuitoare l-a făcut să sufere mult.

%%K t04-c2-10-2 p
Și totuși cu greu s-a înzdrăvenit. Acum îi este mai bine. Dar piciorul, pe care îl
ține pe o \VUgras{pernă} \textsuperscript{(38)} pusă pe un \VUgras{scăunel}
scund \textsuperscript{(39)}, încă îl doare.

%%K t04-c2-11-1 p
11. --- Când este învelit cu grijă în \VUgras{halatul} său cald
\textsuperscript{(35)} cu \VUgras{nasturi} mari de sidef \textsuperscript{(36)}
și când lângă el, ca să îi citească ziarul, stă mica \VUgras{orfană}
\textsuperscript{(40)} cu \VUgras{scufia} albă, cu \VUgras{rochia} albastră
\textsuperscript{(42)} și cu \VUgras{pelerina} \textsuperscript{(41)}, nu se
plânge.

%%K t04-c2-12-1 p
12. --- În fiecare an, când \VUgras{calendarul} \textsuperscript{(43)} aduce din
nou \VUgras{ziua} întâi a lui \VUgras{decembrie}, își așteaptă cu nerăbdare
\VUgras{nepoții}, căci știe că nu vor uita această \VUgras{zi} însemnată.

%%K t04-c2-13-1 p
13. --- Își amintește cu emoție zilele îndepărtate când mergea el însuși să îl
felicite pe vestitul \VUgras{mareșal} al împăratului, al cărui \VUgras{portret}
\textsuperscript{(44)} atârnă pe perete și care le este copiilor
\VUgras{străbunic}.
\VUsaut{6.0mm}
\VUfilet{22mm}
